\documentclass[11pt,a4paper]{article}
\usepackage[utf8]{inputenc}
\usepackage[T1]{fontenc}
\usepackage{amsmath,amssymb,amsthm}
\usepackage{graphicx}
\usepackage{booktabs}
\usepackage{hyperref}
\usepackage[margin=1in]{geometry}
\usepackage{caption}

\newtheorem{theorem}{Theorem}
\newtheorem{proposition}{Proposition}

\title{Optimization-Induced Dynamics and the Inverted-U:\\A Distillation-Based Theory of Optimal Misalignment}
\author{SAI\\{\small Independent Research}}
\date{}

\begin{document}
\maketitle

\begin{abstract}
Conventional approaches to AI alignment aim to minimize divergence between systems,
implicitly assuming that perfect agreement maximizes utility. In this work, we challenge this
assumption by proposing Ma Theory, which posits an inverted-U relationship between structural
misalignment ($\varepsilon^*$) and emergent utility ($E$).

We introduce a Dual-Track Design to isolate the role of generative dynamics: (A) a distillation
track based on stochastic gradient descent (SGD) trajectories across Pythia model checkpoints
(step 0--128,000), and (B) an interpolation track using linear weight interpolation in LoRA
parameter space. Experiments follow a pre-registered protocol with 50 samples across 3 seeds
evaluated on WikiText-2.

We find that the distillation track exhibits a statistically significant inverted-U relationship
in Pythia-410m (post-peak decline: $r = -0.993$, $p = 0.0007$), with a peak at
$\varepsilon^* \approx 68$. The unimodal structure is replicated in Pythia-160m under identical
experimental conditions, with a peak at $\varepsilon^* \approx 35$ (step 32,000), demonstrating
that the phenomenon is replicated across model scales. In contrast, the interpolation track shows
monotonic increase with no evidence of a peak.

We show that this difference is governed by the dynamics of novelty ($D_\mathrm{nov}$):
distillation induces early saturation ($D_\mathrm{nov} \approx 0.98$), while interpolation
produces approximately linear growth without saturation. This saturation is a key enabling
condition for the emergence of the inverted-U.

While multiple asymmetric functions can fit the observed curve, the theoretically derived form
$E = C\varepsilon^*\exp(-k\varepsilon^*)$ provides an interpretable minimal model, whereas
log-normal variants offer higher descriptive flexibility under certain fitting conditions.

We conclude that the inverted-U is not observed along linear interpolation paths in parameter
space, but appears only under optimization-induced dynamics such as SGD. The experimental
evidence is consistent with the choice of generative dynamics playing a causal role in the
formation of the inverted-U structure.
\end{abstract}

%% ============================================================
\section{Introduction}

Recent AI alignment and model integration methods \cite{ouyang2022,christiano2017} have treated
the minimization of distributional divergence---specifically, the KL divergence
($\varepsilon^*$) between teacher and student models---as the ideal objective. However, the Ma
Theory proposed in this paper identifies a \emph{translation paradox}: pursuing ``perfect
agreement'' eliminates emergent utility, which arises from the productive interaction between
systems. We call this framework \emph{Ma Theory}, from the concept of \emph{ma}
(`in-between')---the productive space between cognitive systems from which emergence arises. This study demonstrates the existence of an inverted-U relationship between structural
divergence ($\varepsilon^*$) and emergent utility ($E$) \cite{page2007,wei2022}, where
emergence is maximized at an optimal degree of imperfection.

The central contribution of this paper is the mathematical and experimental demonstration that
this inverted-U relationship is not observed under linear weight interpolation in parameter
space, but appears only through the optimization-induced exploration trajectory of
stochastic gradient descent (SGD) in the distillation process. Weight interpolation corresponds
to a linear path in parameter space, while distillation corresponds to a qualitatively different
exploration trajectory driven by SGD. This contrast enables us to test whether the emergence of
the inverted-U depends on the geometric properties of the learning path. We further replicate
the phenomenon across two model scales (Pythia-410m and Pythia-160m) to assess its generality.
We adopt a preregistered protocol, rigorously evaluating the existence of the inverted-U (H1)
against the null hypothesis of monotonic increase (H0), as well as the dynamical limitations of
the generative process (H2.5), under a falsifiable experimental design.

%% ============================================================
\section{Theory}

This section defines the fundamental theorems constituting Ma Theory and the mathematical model
governing the dynamics of emergent utility.

We define the structural divergence $\varepsilon^*$ between systems $A$ and $B$ as the KL
divergence \cite{cover2006} between their output probability distributions:
\[
\varepsilon^* = D_\mathrm{KL}(p_A \| p_B)
\]

The emergent utility $E$ (specifically, agreement-based emergence $E_\mathrm{agree}$) is
decomposed into the product of embedding-based novelty $D_\mathrm{nov}$ and output token
agreement probability $P_\mathrm{agree}$:
\[
E = D_\mathrm{nov} \times P_\mathrm{agree}
\]

The embedding vector $\phi$ and novelty $D_\mathrm{nov}$ are defined as:
\[
\phi(x;\theta) = \mathrm{L2Norm}\!\left(\frac{1}{T}\sum_{t=1}^{T} h_t^{(L)}(x;\theta)\right),
\quad D_\mathrm{nov} = 1 - \phi_A \cdot \phi_B
\]
where $h_t^{(L)}$ is the hidden state vector at token $t$ in the final hidden layer, and
$\mathrm{L2Norm}$ denotes normalization by the Euclidean norm. The agreement probability
$P_\mathrm{agree}$ is defined as:
\[
P_\mathrm{agree} = \frac{1}{T}\sum_{t=1}^{T}
\mathbf{1}\!\left[\arg\max_v p_A(v|x_{<t}) = \arg\max_v p_B(v|x_{<t})\right]
\]

\begin{theorem}[Complete Translation Impossibility]
When cognitive systems $A$ and $B$ are distinct ($A \neq B$), the structural divergence is
always positive: $\varepsilon^* > 0$.
\end{theorem}

\begin{proof}
By assumption, $A \neq B$ implies $p_A \neq p_B$, i.e., there exists at least one context
$x_{<t}$ and token $v$ such that $p_A(v|x_{<t}) \neq p_B(v|x_{<t})$. By Gibbs' inequality,
$D_\mathrm{KL}(p_A \| p_B) \geq 0$ with equality if and only if $p_A = p_B$. Since
$p_A \neq p_B$, we obtain $\varepsilon^* = D_\mathrm{KL}(p_A \| p_B) > 0$.
\end{proof}

\begin{proposition}[Optimal Imperfection]\label{prop:optimal}
Under distillation dynamics, the emergent utility $E$ forms an inverted-U as a function of
$\varepsilon^*$. The generalized functional form is:
\[
E(\varepsilon^*) = C\,(\varepsilon^*)^\alpha\,\exp(-k\varepsilon^*)
\]
where $C > 0$ is a scaling constant, $k > 0$ governs the decay rate, and $\alpha > 0$
controls the rise exponent. As the theoretical baseline, we fix $\alpha = 1$, yielding the
minimal two-parameter form $E = C\varepsilon^*\exp(-k\varepsilon^*)$ with optimal divergence
$\varepsilon^*_\mathrm{opt} = 1/k$.
\end{proposition}

\noindent\textbf{Note:} This proposition provides a theoretically grounded minimal model
derived from rate-distortion theory \cite{berger1971}, whose parameters ($C$, $k$) admit
physical interpretation and which yields the optimal divergence
$\varepsilon^*_\mathrm{opt} = 1/k$ as a direct prediction. In data-driven fitting, a
log-normal functional form (LogPeak) may yield a lower AIC in some conditions. However,
regardless of the functional form employed, the inverted-U structure decisively rejects
quadratic and linear models. Experimentally, $\alpha$ is not stable across model scales
($\alpha \approx 0.81$ for 410m, $\alpha \approx 1.62$ for 160m under unweighted fitting),
and the AIC ranking of specific functional forms depends on both the fitting method and
model scale. Proposition~\ref{prop:optimal} should therefore be regarded as a theoretically
grounded minimal model rather than a universally optimal functional form.

\begin{theorem}[Non-Understanding as Source of Emergence]\label{thm:nonunderstanding}
Under the assumption that the embedding map $\phi$ is injective with respect to the output
distribution $p$ (i.e., $\phi_A \neq \phi_B$ implies $p_A \neq p_B$), a necessary condition
for the observation of positive emergence is the existence of incomplete understanding
(divergence) between systems: $E > 0 \implies \varepsilon^* > 0$.
\end{theorem}

\begin{proof}
If $E = D_\mathrm{nov} \times P_\mathrm{agree} > 0$, then both factors are positive.
$D_\mathrm{nov} > 0$ implies $\phi_A \cdot \phi_B < 1$, hence $\phi_A \neq \phi_B$.
By the injectivity assumption on $\phi$, distinct embeddings imply distinct output
distributions: $p_A \neq p_B$. By Theorem~1, $\varepsilon^* > 0$.

\medskip
\noindent\textbf{Remark on the injectivity assumption.}
In the architectures considered (Pythia/GPT-NeoX), the output distribution is computed as
$p = \mathrm{softmax}(W_\mathrm{unembed}\, h^{(L)})$, and the embedding $\phi$ is a
deterministic function of $h^{(L)}$ involving token-wise averaging and L2 normalization.
The natural direction of dependence is $h^{(L)} \to p$ and $h^{(L)} \to \phi$. The injectivity assumption ($\phi_A \neq \phi_B \Rightarrow p_A \neq p_B$) is adopted as an
architecture-dependent working assumption in this paper. Because $\phi$ involves averaging
over token positions, information may in principle be lost, and a full-rank unembedding
matrix alone does not guarantee injectivity of the composed map. Empirical verification or
the identification of stronger sufficient conditions is left to future work.
\end{proof}

\paragraph{Dynamical Derivation of the Inverted-U (The Friction Update).}
The conditions for the emergence of the inverted-U described in
Proposition~\ref{prop:optimal} are mathematically explained by the following differential
equation. We model the discrete checkpoints as samples from a continuous, differentiable
learning trajectory parameterized by $\varepsilon^*$:
\[
\frac{dE}{d\varepsilon^*} =
\frac{dD_\mathrm{nov}}{d\varepsilon^*} \times P_\mathrm{agree}
+ D_\mathrm{nov} \times \frac{dP_\mathrm{agree}}{d\varepsilon^*}
\]

In distillation dynamics (SGD trajectory), the novelty $D_\mathrm{nov}$ saturates at its
upper bound ($\approx 0.98$) during early learning, yielding
$\frac{dD_\mathrm{nov}}{d\varepsilon^*} \approx 0$. Consequently, the second term (negative
contribution) becomes dominant, forming the right-side decay region of the inverted-U.

In contrast, under weight interpolation (linear path in parameter space),
$D_\mathrm{nov} \propto \varepsilon^*$ (linear increase) without saturation. The positive
contribution of the first term consistently dominates, resulting in
$\frac{dE}{d\varepsilon^*} > 0$ (perpetual monotonic increase). Therefore, the inverted-U is
not expected to form along linear interpolation paths in parameter space, and is not observed
in our experiments.

%% ============================================================
\section{Methods}

\subsection{Dual-Track Experimental Design}

This study adopts a ``Dual-Track Design'' that examines the relationship between structural
divergence $\varepsilon^*$ and emergent utility $E$ under two fundamentally different
generative processes.

\textbf{Track A: Distillation} measures the divergence $\varepsilon^*$ between each
checkpoint and the final model as a single neural network progresses from initialization to
its final state through pretraining. This trajectory is an optimization-induced exploration
process driven by SGD \cite{hinton2015}.

\textbf{Track B: Interpolation} controls $\varepsilon^*$ by linearly interpolating the
weights of a base model and a LoRA fine-tuned model, varying the interpolation coefficient
$\lambda \in [0,1]$. Linear movement in parameter space constitutes a Euclidean geodesic
in parameter space \cite{amari2016,frankle2020,wortsman2022}, though it does not in general
correspond to a geodesic on the statistical manifold under Fisher information geometry
\cite{amari2000}.

This contrast tests whether the emergence of the inverted-U depends not only on the
magnitude of $\varepsilon^*$ but on the geometric properties of the path that generates it.

\subsection{Models}

All experiments use EleutherAI/Pythia models \cite{biderman2023}, autoregressive language
models based on the GPT-NeoX architecture \cite{black2022}. All pretraining checkpoints
(step 0--143,000) are publicly available on HuggingFace.

\textbf{Primary scale (410m):} For Track A, step 143,000 serves as the teacher model, with
steps \{0, 1000, 2000, 4000, 8000, 16000, 32000, 64000, 128000\} as student models. Step 0
corresponds to the fully initialized state (random weights). For Track B, a student model was
created by applying LoRA fine-tuning \cite{hu2022} (rank=64, $\alpha$=256, target modules:
query, key, value) to the Pythia-410m base model. The interpolated model weights are
constructed as $W(\lambda) = W_\mathrm{base} + \lambda \cdot \Delta W$, with 16 values of
$\lambda$ measured.

LoRA fine-tuning was conducted on \texttt{code\_search\_net/python} (train split, first
5,000 samples) with a constant learning rate of $1 \times 10^{-4}$ (no scheduler), a custom
training loop (manual backward pass), for 10,000 steps. A preliminary verification confirmed
that \texttt{merge\_and\_unload()} has negligible effect on $\varepsilon^*$
(merge/LoRA ratio = 0.9997).

\textbf{Replication scale (160m):} To assess scale generality, the distillation track was
replicated using Pythia-160m with identical protocol: teacher at step 143,000, the same 9
student checkpoints, and identical evaluation conditions (WikiText-2 test, 50 samples
$\times$ 3 seeds).

\subsection{Measurement Protocol}

All experiments use a unified measurement codebase. Evaluation data consists of sequences
extracted from the test split of WikiText-2. The following four metrics are measured:

\textbf{Structural divergence} $\varepsilon^*$:
\[
\varepsilon^* = \sum_{t=1}^{T} D_\mathrm{KL}(p_T(\cdot|x_{<t}) \| p_S(\cdot|x_{<t}))
\]

\textbf{Novelty} $D_\mathrm{nov}$:
\[
\phi(x;\theta) = \mathrm{L2Norm}\!\left(\frac{1}{T}\sum_{t=1}^{T}
h_t^{(L)}(x;\theta)\right), \quad D_\mathrm{nov} = 1 - \phi_T \cdot \phi_S
\]

\textbf{Agreement probability} $P_\mathrm{agree}$:
\[
P_\mathrm{agree} = \frac{1}{T}\sum_{t=1}^{T}
\mathbf{1}\!\left[\arg\max_v p_T(v|x_{<t}) = \arg\max_v p_S(v|x_{<t})\right]
\]

\textbf{Emergent utility} $E_\mathrm{agree}$:
\[
E_\mathrm{agree} = D_\mathrm{nov} \times P_\mathrm{agree}
\]

All four metrics are computed per input sequence $x$. The reported values are sample means
over $N=50$ sequences drawn from the evaluation dataset $\mathcal{D}$:
\[
\bar{E}_\mathrm{agree} = \frac{1}{N}\sum_{i=1}^{N} E_\mathrm{agree}(x_i),
\quad x_i \sim \mathcal{D}
\]
where $E_\mathrm{agree}(x_i) = D_\mathrm{nov}(x_i) \times P_\mathrm{agree}(x_i)$ is
computed independently for each sequence. This procedure is repeated for each of 3 seeds
(controlling sequence sampling), and the final estimates are seed-averaged with standard
errors.

Each condition is measured with 50 samples $\times$ 3 seeds, reporting means and standard
errors. KL divergence is computed in float32 precision. Each input sequence is truncated to a
maximum of 512 tokens without padding; sequences shorter than 512 tokens retain their natural
length. Since $\varepsilon^*$ is defined as a sum over tokens, it depends on the effective
sequence length $T$. However, because all metrics are computed per sequence and then averaged
across sequences, the per-sample variation in $T$ is absorbed into the sample mean, and
cross-condition comparisons (which use identical input sequences) are not affected.

\subsection{Pre-registration and Hypothesis Testing}

To eliminate cherry-picking, the following hypothesis testing structure was preregistered:

\begin{itemize}
\item \textbf{H0 (Null):} $E_\mathrm{agree}$ increases monotonically with $\varepsilon^*$;
  no inverted-U exists.
\item \textbf{H1a:} $E_\mathrm{agree}$ exhibits an inverted-U with an interior peak
  $\varepsilon^*_\mathrm{opt}$.
\item \textbf{H1b:} The functional form of Proposition~\ref{prop:optimal} provides the
  best fit among candidate models.
\item \textbf{H2.5:} The emergence of the inverted-U depends on the generative process
  (distillation vs.\ interpolation).
\end{itemize}

Rejection of H0 uses a Pearson correlation significance test on the post-peak decay, with
Welch's $t$-test for adjacent-point comparisons.

\subsection{Model Comparison}

To evaluate the functional form of Proposition~\ref{prop:optimal}, five candidate models are
compared using the Akaike Information Criterion (AIC) \cite{akaike1974}:

\begin{enumerate}
\item Prop.~\ref{prop:optimal} original (2p): $E = C\varepsilon^*\exp(-k\varepsilon^*)$
\item Prop.~\ref{prop:optimal} gamma (3p): $E = C(\varepsilon^*)^\alpha\exp(-k\varepsilon^*)$
\item Log-peak (3p): $E = C\exp\!\left(-\frac{(\ln\varepsilon^* - \mu)^2}{2\sigma^2}\right)$
\item Quadratic (3p): $E = a(\varepsilon^*)^2 + b\varepsilon^* + c$
\item Linear (2p): $E = a\varepsilon^* + b$
\end{enumerate}

Model comparison is conducted under both weighted least squares (inverse standard error
weighting) and unweighted least squares to verify the stability of model rankings across
fitting methods.

%% ============================================================
\section{Results}

\subsection{Distillation Track: Confirmation of the Inverted-U (Pythia-410m)}

Figure~\ref{fig:410m} and Table~\ref{tab:410m} present all measurements from the distillation
track. $E_\mathrm{agree}$ exhibits a clear inverted-U with respect to $\varepsilon^*$, peaking
at step 8,000 ($\varepsilon^* = 67.57$, $E_\mathrm{agree} = 0.468$).

\textbf{Statistical verification:} The correlation between $E_\mathrm{agree}$ and
$\varepsilon^*$ in the post-peak region is $r = -0.993$, $p = 0.0007$, confirming a highly
significant decay trend. Adjacent-point differences near the peak do not reach significance
(peak vs.\ step 16,000: $t = 0.843$, $p = 0.20$; peak vs.\ step 4,000: $t = 0.394$,
$p = 0.35$), consistent with a broad plateau in the range
$\varepsilon^* \approx 50$--$90$.

The null hypothesis H0 (monotonic increase) is rejected at $p = 0.0007$.

\begin{figure}[t]
\centering
\includegraphics[width=0.85\textwidth]{fig1_410m_eagree.pdf}
\caption{Pythia-410m distillation track: $E_\mathrm{agree}$ vs.\ $\varepsilon^*$.
The inverted-U peaks at step 8,000 ($\varepsilon^* = 67.57$). Step~0
($\varepsilon^* = 969.9$, $E_\mathrm{agree} = 0$) is omitted for readability;
see Table~\ref{tab:410m}. Error bars indicate standard errors of $\varepsilon^*$
across 3 seeds.}
\label{fig:410m}
\end{figure}

\begin{table}[t]
\centering
\caption{Distillation track (Pythia-410m): All measurements.
$\star$ marks the peak of $E_\mathrm{agree}$.}
\label{tab:410m}
\begin{tabular}{rcccc}
\toprule
Step & $\varepsilon^*$ (mean $\pm$ SE) & $D_\mathrm{nov}$ & $P_\mathrm{agree}$
& $E_\mathrm{agree}$ \\
\midrule
0       & $969.90 \pm 4.01$   & 0.982 & 0.000 & 0.000 \\
1,000   & $248.67 \pm 2.73$   & 0.886 & 0.366 & 0.325 \\
2,000   & $140.11 \pm 1.73$   & 0.828 & 0.497 & 0.411 \\
4,000   & $92.60 \pm 1.16$    & 0.815 & 0.571 & 0.465 \\
8,000   & $67.57 \pm 0.86$    & 0.745 & 0.628 & $0.468^\star$ \\
16,000  & $49.91 \pm 0.68$    & 0.684 & 0.674 & 0.462 \\
32,000  & $38.77 \pm 0.56$    & 0.638 & 0.713 & 0.455 \\
64,000  & $23.97 \pm 0.48$    & 0.458 & 0.765 & 0.351 \\
128,000 & $5.76 \pm 0.09$     & 0.012 & 0.874 & 0.010 \\
\bottomrule
\end{tabular}
\end{table}

\subsection{Scaling Replication (Pythia-160m)}

To assess whether the inverted-U structure generalizes beyond a single model scale, the
distillation track was replicated using Pythia-160m under identical experimental conditions
(Figure~\ref{fig:160m} and Table~\ref{tab:160m}).

The inverted-U structure is replicated: $E_\mathrm{agree}$ peaks at step 32,000
($\varepsilon^* = 34.96$, $E_\mathrm{agree} = 0.644$), with a post-peak monotonic decline
observed across the subsequent two intervals. Minor local fluctuations are present in the
pre-peak region (a slight decrease from step 1,000 to step 2,000), but do not affect the
overall unimodal structure.

A key observation is that $D_\mathrm{nov}$ remains in a high-value regime ($> 0.90$) from
step 1,000 through step 32,000 in the 160m model, compared to rapid decay beginning at
step 4,000 in the 410m model. This extended saturation is consistent with a delayed
transition to negative-dominated dynamics, corresponding to a later peak. The mechanism is
fully consistent with the Friction Update dynamics derived in Section~2.

\begin{figure}[t]
\centering
\includegraphics[width=0.85\textwidth]{fig2_160m_eagree.pdf}
\caption{Pythia-160m distillation track: $E_\mathrm{agree}$ vs.\ $\varepsilon^*$.
The inverted-U is replicated, peaking at step 32,000 ($\varepsilon^* = 34.96$).
Step~0 ($\varepsilon^* = 690.2$, $E_\mathrm{agree} = 0$) is omitted for readability;
see Table~\ref{tab:160m}.}
\label{fig:160m}
\end{figure}

\begin{table}[t]
\centering
\caption{Distillation track (Pythia-160m): All measurements.
$\star$ marks the peak of $E_\mathrm{agree}$.}
\label{tab:160m}
\begin{tabular}{rcccc}
\toprule
Step & $\varepsilon^*$ (mean $\pm$ SE) & $D_\mathrm{nov}$ & $P_\mathrm{agree}$
& $E_\mathrm{agree}$ \\
\midrule
0       & $690.21 \pm 112.35$ & 0.987 & 0.000 & 0.000 \\
1,000   & $128.02 \pm 20.28$  & 0.919 & 0.453 & 0.415 \\
2,000   & $77.65 \pm 11.71$   & 0.910 & 0.443 & 0.400 \\
4,000   & $54.40 \pm 8.31$    & 0.934 & 0.580 & 0.541 \\
8,000   & $46.17 \pm 7.00$    & 0.984 & 0.627 & 0.617 \\
16,000  & $39.61 \pm 6.00$    & 0.991 & 0.639 & 0.633 \\
32,000  & $34.96 \pm 5.41$    & 0.971 & 0.664 & $0.644^\star$ \\
64,000  & $27.46 \pm 4.27$    & 0.904 & 0.611 & 0.549 \\
128,000 & $12.03 \pm 1.84$    & 0.002 & 0.725 & 0.001 \\
\bottomrule
\end{tabular}
\end{table}

\subsection{Cross-Scale Comparison}

Table~\ref{tab:crossscale} compares the distillation track results across the two model scales.

The unimodal structure of $E_\mathrm{agree}$ is replicated across model scales under
identical experimental conditions, confirming that the phenomenon is not an artifact of a
particular model size. However, two scale-dependent differences emerge:

\textbf{Peak shift.} The peak occurs at step 32,000 ($\varepsilon^* = 34.96$) for 160m
versus step 8,000 ($\varepsilon^* = 67.57$) for 410m. The present results suggest that
$\varepsilon^*_\mathrm{opt}$ and the corresponding training step may depend on model scale,
though this observation is based on two scales and should be regarded as preliminary.

\textbf{$D_\mathrm{nov}$ saturation range.} The high-$D_\mathrm{nov}$ regime ($> 0.9$)
persists from step 0 to step 32,000 in the 160m model, compared to step 0 to step 4,000 in
the 410m model. This difference in saturation range directly governs the peak shift through
the Friction Update dynamics: as long as $D_\mathrm{nov}$ remains saturated
($\frac{dD_\mathrm{nov}}{d\varepsilon^*} \approx 0$), $E_\mathrm{agree}$ continues to rise
with increasing $P_\mathrm{agree}$.

\textbf{Statistical note.} For the 160m model, the post-peak trajectory exhibits a strict
monotonic decrease across the final two intervals. While the delayed peak structurally limits
the number of post-peak data points ($n = 3$ including the peak), the observed downward trend
is consistent with the 410m results, where four post-peak points yield $r = -0.993$,
$p = 0.0007$.

\begin{table}[t]
\centering
\caption{Cross-scale comparison of distillation track results.}
\label{tab:crossscale}
\begin{tabular}{lcc}
\toprule
& Pythia-160m & Pythia-410m \\
\midrule
Peak step             & 32,000  & 8,000  \\
$\varepsilon^*$ at peak & 34.96 & 67.57  \\
$E_\mathrm{agree}$ at peak & 0.644 & 0.468 \\
$\varepsilon^*$ range & 12.0--690 & 5.8--970 \\
Post-peak monotonic decline & Yes & Yes \\
Post-peak data points & 2 & 4 \\
$D_\mathrm{nov} > 0.9$ range & steps 0--32k & steps 0--4k \\
\bottomrule
\end{tabular}
\end{table}

\subsection{Interpolation Track: Monotonic Increase (Negative Control)}

Figure~\ref{fig:interp} and Table~\ref{tab:interp} present measurements from the
interpolation track (Track B, Pythia-410m). $E_\mathrm{agree}$ increases monotonically with
$\varepsilon^*$, with no peak observed up to $\varepsilon^* = 159.44$.

This result demonstrates that the emergence of the inverted-U depends on the generative
process, supporting hypothesis H2.5 (process dependence).

\begin{figure}[t]
\centering
\includegraphics[width=0.85\textwidth]{fig3_interp_eagree.pdf}
\caption{Interpolation track (Pythia-410m): $E_\mathrm{agree}$ vs.\ $\varepsilon^*$.
Monotonic increase with no peak---negative control confirming process dependence.}
\label{fig:interp}
\end{figure}

\begin{table}[t]
\centering
\caption{Interpolation track (Pythia-410m): Selected measurements.}
\label{tab:interp}
\begin{tabular}{ccccc}
\toprule
$\lambda$ & $\varepsilon^*$ & $D_\mathrm{nov}$ & $P_\mathrm{agree}$
& $E_\mathrm{agree}$ \\
\midrule
0.10 & 11.75  & 0.025 & 0.838 & 0.021 \\
0.30 & 25.49  & 0.048 & 0.762 & 0.036 \\
0.50 & 45.51  & 0.072 & 0.695 & 0.050 \\
0.70 & 77.21  & 0.105 & 0.623 & 0.066 \\
1.00 & 159.44 & 0.190 & 0.494 & 0.094 \\
\bottomrule
\end{tabular}
\end{table}

\subsection{$D_\mathrm{nov}$ Saturation Dynamics: Mechanism of the Inverted-U}

The behavior of $D_\mathrm{nov}$ differs qualitatively between the distillation and
interpolation tracks (Figure~\ref{fig:dnov} and Table~\ref{tab:dnov}).

In the distillation track, $D_\mathrm{nov}$ reaches near its upper bound (0.982) at step 0
($\varepsilon^* = 969.9$) and decays as learning progresses, exhibiting early saturation. In
the interpolation track, $D_\mathrm{nov} \propto \varepsilon^*$ holds as a linear
relationship, with no saturation observed.

As derived in the Theory section, when $D_\mathrm{nov}$ saturates
($\frac{dD_\mathrm{nov}}{d\varepsilon^*} \approx 0$), the negative second term of the
dynamical equation dominates, forming the right-side decay of the inverted-U. When
$D_\mathrm{nov}$ does not saturate (interpolation), the positive first term consistently
prevails, maintaining monotonic increase.

\begin{figure}[t]
\centering
\includegraphics[width=0.85\textwidth]{fig4_dnov_comparison.pdf}
\caption{$D_\mathrm{nov}$ vs.\ $\varepsilon^*$: distillation (410m, 160m) vs.\
interpolation (410m). Distillation exhibits early saturation ($D_\mathrm{nov} > 0.9$),
while interpolation shows linear growth without saturation. The dashed line marks
$D_\mathrm{nov} = 0.9$.}
\label{fig:dnov}
\end{figure}

\begin{table}[t]
\centering
\caption{Comparison of $D_\mathrm{nov}$ at matched $\varepsilon^*$ regions
(distillation vs.\ interpolation, 410m).}
\label{tab:dnov}
\begin{tabular}{cccc}
\toprule
$\varepsilon^*$ & $D_\mathrm{nov}$ (distill.) & $D_\mathrm{nov}$ (interp.) & Ratio \\
\midrule
$\sim 7$   & 0.012 & 0.014 & $0.8\times$ \\
$\sim 25$  & 0.458 & 0.048 & $9.6\times$ \\
$\sim 45$  & 0.684 & 0.072 & $9.6\times$ \\
$\sim 77$  & 0.745 & 0.105 & $7.1\times$ \\
\bottomrule
\end{tabular}
\end{table}

\subsection{Model Comparison}

Table~\ref{tab:aic} presents the AIC comparison under both fitting conditions for the 410m
data (steps 1,000--128,000, step 0 excluded).

Under weighted fitting, Proposition~\ref{prop:optimal} gamma (3p) achieves the best AIC; under unweighted
fitting, LogPeak (3p) achieves the best AIC. The model ranking thus depends on the fitting
method. However, under both conditions, all unimodal asymmetric models decisively reject
quadratic and linear models ($\Delta\mathrm{AIC} > 20$ under unweighted conditions),
confirming that the inverted-U structure is robust to model selection.

For the 160m data, LogPeak (3p) achieves the best AIC under both fitting conditions
($R^2 = 0.93$ unweighted, $R^2 = 0.99$ weighted), with Proposition~\ref{prop:optimal} models showing
$\Delta\mathrm{AIC} > 11$. This indicates that the AIC ranking of specific functional forms
is not stable across model scales, reinforcing the position that Proposition~\ref{prop:optimal}
should be regarded as a theoretically grounded minimal model rather than a universally optimal
functional form.

\begin{table}[t]
\centering
\caption{AIC model comparison (Pythia-410m). $\Delta$AIC is computed from the best model
within each condition.}
\label{tab:aic}
\begin{tabular}{lccccc}
\toprule
& \multicolumn{2}{c}{Weighted (inv.\ SE)} & & \multicolumn{2}{c}{Unweighted} \\
\cmidrule{2-3}\cmidrule{5-6}
Model & AIC & $\Delta$AIC & & AIC & $\Delta$AIC \\
\midrule
Prop.~\ref{prop:optimal} gamma (3p) & 16.12 & 0.00  & & $-41.17$ & 11.64 \\
Prop.~\ref{prop:optimal} orig (2p)  & 32.73 & 16.61 & & $-41.84$ & 10.96 \\
LogPeak (3p)       & 28.71 & 12.59 & & $-52.81$ & 0.00  \\
Quadratic (3p)     & 27.37 & 11.25 & & $-31.35$ & 21.45 \\
Linear (2p)        & 33.52 & 17.40 & & $-27.21$ & 25.59 \\
\bottomrule
\end{tabular}
\end{table}

\subsection{Stability of $\alpha$}

The estimated $\alpha$ in the generalized Proposition~\ref{prop:optimal} gamma shows sensitivity to both the
fitting method and model scale (Table~\ref{tab:alpha}).

Under unweighted fitting for the 410m data, $\alpha$ is stable at $0.81 \pm 0.03$ across all
six robustness conditions (sample sizes $n = 25, 50, 100$; sequence lengths $T = 64, 128,
256$). The AIC difference between Proposition~\ref{prop:optimal} original ($\alpha = 1$ fixed, 2p) and
Proposition~\ref{prop:optimal} gamma ($\alpha$ free, 3p) is $\Delta\mathrm{AIC} = 0.67$ (negligible),
indicating that $\alpha = 1$ provides a sufficient approximation at the 410m scale.

However, at the 160m scale, the unweighted estimate yields $\alpha = 1.62$, indicating that
$\alpha$ is not stable across model scales. This scale dependence further supports treating
Proposition~\ref{prop:optimal} as a useful theoretical framework rather than a universal
empirical law.

\begin{table}[t]
\centering
\caption{Estimated $\alpha$ under different fitting conditions and model scales.}
\label{tab:alpha}
\begin{tabular}{lcc}
\toprule
Fitting condition & $\alpha$ (410m) & $\alpha$ (160m) \\
\midrule
Weighted (incl.\ step 0) & 2.34 & --- \\
Weighted (excl.\ step 0) & 2.12 & $12.50^\dagger$ \\
Unweighted (all conditions) & $0.81 \pm 0.03$ & 1.62 \\
\bottomrule
\multicolumn{3}{l}{\footnotesize $^\dagger$Degenerate fit ($C \approx 0$); not interpretable.}
\end{tabular}
\end{table}

\subsection{Additional Findings}

\textbf{$\varepsilon^*(\lambda) \propto \lambda^2$:} In the interpolation track,
$\varepsilon^*$ follows a quadratic function of $\lambda$ with high precision (Quadratic
$R^2 = 0.9998$, Linear $R^2 = 0.9609$), consistent with the Fisher information matrix
approximation \cite{amari2000}
$\varepsilon^* \approx \frac{1}{2}\Delta W^T G(\theta)\Delta W$. This result was consistently
confirmed across five interpolation experiments with different LoRA configurations. Notably,
this quadratic relationship confirms that linear interpolation in parameter space does not
correspond to a geodesic on the statistical manifold (which would yield
$\varepsilon^* \propto \lambda$), but rather traces a nonlinear path in distribution space.

\textbf{Effect of training conditions on $\varepsilon^*$:} Under identical LoRA settings
(rank=64, $\alpha$=256, 10,000 steps), changing training conditions (learning rate, scheduler,
training data size) altered $\varepsilon^*$ at $\lambda = 1.0$ by a factor of 6.8 (from 23.35
to 159.44), while the weight norm $\|\Delta W\|$ changed by only $1.15\times$. Thus,
$\varepsilon^*$ depends on the direction rather than the magnitude of $\Delta W$.

\textbf{Domain independence:} Evaluating a CODE-trained LoRA model on both CODE and WikiText
yielded a scale ratio of $0.9\times$, indicating that evaluation domain has minimal effect on
$\varepsilon^*$.

\textbf{Seed independence:} No significant seed effects were detected in any experiment
(ANOVA $p > 0.05$ for all metrics across all conditions).

\subsection{Hypothesis Verdict}

\begin{table}[t]
\centering
\caption{Preregistered hypothesis verdicts.}
\label{tab:verdict}
\begin{tabular}{lll}
\toprule
Hypothesis & Verdict & Evidence \\
\midrule
H0 (monotonic) & Rejected & Post-peak $r = -0.993$, $p = 0.0007$ (410m) \\
H1a (inverted-U exists) & Confirmed & Inverted-U confirmed at two scales \\
H1b (Prop.~\ref{prop:optimal} form) & Not established & AIC rank depends on method and scale \\
H2.5 (process dep.) & Confirmed & Interpolation: monotonic; Distillation: inverted-U \\
\bottomrule
\end{tabular}
\end{table}

The existence of the inverted-U structure (H1a) is confirmed across two model scales. The
specific functional form of Proposition~\ref{prop:optimal} (H1b) is not established as
universally optimal; its AIC ranking varies with fitting method and model scale.

%% ============================================================
\section{Discussion}

This study rejected the conventional hypothesis (H0) that complete agreement
($\varepsilon^* \to 0$) maximizes emergence, and statistically confirmed the inverted-U
phenomenon in which emergence is maximized at optimal imperfection ($\varepsilon^* > 0$). The
phenomenon is replicated across two model scales, strengthening the case for its generality.
This section discusses the dynamical conditions producing this phenomenon and the methodological
insights from model selection.

\subsection{Saturation Dynamics of $D_\mathrm{nov}$ and Information-Geometric Interpretation}

The conditions for the inverted-U are unified by the differential equation
$\frac{dE}{d\varepsilon^*} = \frac{dD_\mathrm{nov}}{d\varepsilon^*} \times P_\mathrm{agree}
+ D_\mathrm{nov} \times \frac{dP_\mathrm{agree}}{d\varepsilon^*}$.

Weight interpolation constitutes linear movement in parameter space. Along this trajectory,
representational distance ($D_\mathrm{nov}$) increases linearly with distributional distance
($\varepsilon^*$), so $D_\mathrm{nov}$ does not saturate and
$\frac{dD_\mathrm{nov}}{d\varepsilon^*} > 0$ holds throughout. Consequently, the positive
first-term contribution dominates, and emergence increases monotonically. The proposed mechanism
therefore predicts that the inverted-U does not form along linear interpolation paths, which is
consistent with our experimental observations.

In contrast, distillation is an optimization-induced trajectory exploring the loss landscape
via SGD.
Along this trajectory, the representation space rapidly orthogonalizes during early learning,
and $D_\mathrm{nov}$ saturates at its upper bound. This yields
$\frac{dD_\mathrm{nov}}{d\varepsilon^*} \approx 0$, allowing the negative second term to
dominate and forming the right tail (decay region) of the inverted-U.

The cross-scale comparison reveals that the duration of $D_\mathrm{nov}$ saturation varies
with model scale: the high-$D_\mathrm{nov}$ regime ($> 0.9$) extends from step 0 to step
32,000 in the 160m model but only to step 4,000 in the 410m model. This difference governs
the peak shift through the same dynamical mechanism, providing further evidence consistent with
the Friction Update equation.

\subsection{Model Selection and Methodological Insights on $\alpha$}

AIC-based model comparison revealed that the optimal functional form changes depending on both
the fitting method and the model scale. For the 410m data, weighted least squares selects
Proposition~\ref{prop:optimal}, while unweighted least squares selects LogPeak. For the 160m data, LogPeak
achieves the best AIC under both conditions.

LogPeak provides the highest data-driven fit but lacks theoretical grounding in its parameters.
In contrast, Proposition~\ref{prop:optimal} is a minimal parametric model with a theoretical
foundation in rate-distortion theory \cite{berger1971,shannon1948}. The optimization parameter
$\alpha$ is stable at $\alpha \approx 0.8$ under unweighted conditions for the 410m scale, and
$\Delta\mathrm{AIC} = 0.67$ (negligible) when fixed at $\alpha = 1$. However, $\alpha$ varies
across scales (0.81 for 410m, 1.62 for 160m), indicating that it is not a universal constant.

Crucially, regardless of whether LogPeak or Proposition~\ref{prop:optimal} is employed, and regardless of model
scale, quadratic and linear models are decisively rejected ($\Delta\mathrm{AIC} > 10$). This
highlights a separation between descriptive accuracy and theoretical interpretability. The
emergence function in this system belongs to the class of ``unimodal, asymmetric,
right-tail-decaying functions,'' and the existence of the inverted-U phenomenon is a robust
fact independent of model selection or scale.

\subsection{Limitations}

First, although the inverted-U is confirmed across two model scales (160m and 410m) within the
Pythia family, all experiments use a single architecture (GPT-NeoX). Generalization to
different architectures (e.g., encoder-decoder models, mixture-of-experts) remains unverified.

Second, the ``student models'' in the distillation track represent different learning stages of
the same model, rather than truly independent pretrained systems. Evaluation of dynamics between
independently trained models is a task for future work.

Third, the definition of emergent utility $E = D_\mathrm{nov} \times P_\mathrm{agree}$ depends
on specific operational choices. Whether structural properties are preserved under alternative
definitions of $D_\mathrm{nov}$ (e.g., using intermediate-layer representations) or alternative
agreement metrics (e.g., top-$k$ agreement, JSD, cosine similarity between softmax outputs)
requires additional verification.

Fourth, extension of this theory to heterogeneous cognitive systems such as human--AI
collaboration remains at the level of suggestion and requires empirical confirmation.

Fifth, the 160m distillation track provides only two post-peak data points (three including the
peak), limiting the statistical strength of the post-peak decline assessment at this scale. The
observed monotonic decrease is consistent with the 410m results, but additional intermediate
checkpoints would strengthen the evidence.

Sixth, Theorem~\ref{thm:nonunderstanding} relies on the injectivity of the embedding map
$\phi$ with respect to the output distribution $p$. This is adopted as an
architecture-dependent working assumption; because $\phi$ involves token-wise averaging and
normalization, injectivity is not guaranteed in general. The theorem's scope is therefore
limited to architectures and conditions under which this assumption holds.

%% ============================================================
\section{Conclusion}

This study demonstrated the existence of an inverted-U relationship between structural
divergence ($\varepsilon^*$) and emergent utility ($E$) with statistical significance
($p = 0.0007$ for Pythia-410m). The unimodal structure is replicated in Pythia-160m under
identical experimental conditions, supporting the presence of the phenomenon across multiple
model scales. This result shows that emergence depends not on the magnitude of divergence, but
on the nature of the generative dynamics. The experimental evidence is consistent with
generative dynamics playing a causal role in the formation of the inverted-U structure.

The inverted-U appears only under optimization-induced dynamics based on SGD (distillation);
under linear weight interpolation in parameter space, it is not observed.
The dynamical prerequisite is the early saturation of novelty $D_\mathrm{nov}$ accompanying
SGD-driven exploration, as established through the differential equation analysis. The
cross-scale comparison further reveals that the range over which $D_\mathrm{nov}$ remains
saturated varies with model size, governing the position of the optimal divergence
$\varepsilon^*_\mathrm{opt}$ through the same dynamical mechanism.

While multiple asymmetric functions can approximate the phenomenon in a data-driven manner,
Proposition~\ref{prop:optimal} ($E = C\varepsilon^*\exp(-k\varepsilon^*)$) serves as a useful
minimal model with theoretical grounding. Its AIC ranking varies with fitting method and model
scale, and the phenomenon itself---the inverted-U structure and its dependence on
optimization-induced dynamics---is the robust, scale-invariant finding of this work.

%% ============================================================
\section*{Acknowledgements}

This work was developed through a collaborative cognitive framework between a human researcher
and multiple large language model agents (Claude by Anthropic, GPT by OpenAI, and Gemini by
Google). All three AI systems contributed to theoretical development, experimental design,
mathematical verification, and iterative peer review. The author name SAI represents this
collaborative entity.

%% ============================================================
\bibliographystyle{plain}
\begin{thebibliography}{16}

\bibitem{akaike1974}
H.~Akaike.
\newblock A new look at the statistical model identification.
\newblock {\em IEEE Trans.\ Automatic Control}, 19(6):716--723, 1974.

\bibitem{amari2016}
S.~Amari.
\newblock {\em Information Geometry and Its Applications}.
\newblock Springer, 2016.

\bibitem{amari2000}
S.~Amari and H.~Nagaoka.
\newblock {\em Methods of Information Geometry}.
\newblock American Mathematical Society, 2000.

\bibitem{berger1971}
T.~Berger.
\newblock {\em Rate Distortion Theory: A Mathematical Basis for Data
  Compression}.
\newblock Prentice-Hall, 1971.

\bibitem{biderman2023}
S.~Biderman, H.~Schoelkopf, Q.~Anthony, et~al.
\newblock Pythia: A suite for analyzing large language models across training
  and scaling.
\newblock In {\em Proceedings of ICML}, 2023.

\bibitem{black2022}
S.~Black, S.~Biderman, E.~Hallahan, et~al.
\newblock {GPT-NeoX-20B}: An open-source autoregressive language model.
\newblock In {\em ACL Workshop on Challenges in Creating Large Language Models},
  2022.

\bibitem{christiano2017}
P.~F. Christiano, J.~Leike, T.~Brown, et~al.
\newblock Deep reinforcement learning from human preferences.
\newblock {\em Advances in NeurIPS}, 30, 2017.

\bibitem{cover2006}
T.~M. Cover and J.~A. Thomas.
\newblock {\em Elements of Information Theory}.
\newblock Wiley-Interscience, 2nd edition, 2006.

\bibitem{frankle2020}
J.~Frankle, G.~K. Dziugaite, D.~M. Roy, and M.~Carbin.
\newblock Linear mode connectivity and the lottery ticket hypothesis.
\newblock In {\em Proceedings of ICML}, 2020.

\bibitem{hinton2015}
G.~Hinton, O.~Vinyals, and J.~Dean.
\newblock Distilling the knowledge in a neural network.
\newblock {\em arXiv preprint arXiv:1503.02531}, 2015.

\bibitem{hu2022}
E.~J. Hu, Y.~Shen, P.~Wallis, et~al.
\newblock {LoRA}: Low-rank adaptation of large language models.
\newblock In {\em Proceedings of ICLR}, 2022.

\bibitem{ouyang2022}
L.~Ouyang, J.~Wu, X.~Jiang, et~al.
\newblock Training language models to follow instructions with human feedback.
\newblock {\em Advances in NeurIPS}, 35, 2022.

\bibitem{page2007}
S.~E. Page.
\newblock {\em The Difference: How the Power of Diversity Creates Better
  Groups, Firms, Schools, and Societies}.
\newblock Princeton University Press, 2007.

\bibitem{shannon1948}
C.~E. Shannon.
\newblock A mathematical theory of communication.
\newblock {\em Bell System Technical Journal}, 27(3):379--423, 1948.

\bibitem{wei2022}
J.~Wei, Y.~Tay, R.~Bommasani, et~al.
\newblock Emergent abilities of large language models.
\newblock {\em Transactions on Machine Learning Research}, 2022.

\bibitem{wortsman2022}
M.~Wortsman, G.~Ilharco, S.~Y. Gadre, et~al.
\newblock Model soups: Averaging weights of multiple fine-tuned models improves
  accuracy without increasing inference time.
\newblock In {\em Proceedings of ICML}, 2022.

\end{thebibliography}

\end{document}
